\chapter{Discussion and Conclusion}
\label{ch:conclusion}
\iffalse
\fi
\section{Discussion of Results}
\label{sec:5.1}

The four research questions formulated in Chapter~\ref{ch:introduction} guided both the design of the measurement campaign (Chapter~\ref{ch:method}) and the presentation of empirical results (Chapter~\ref{ch:results}). This section interprets those results in light of the existing literature and draws out practical implications for measurement methodology. The campaign spans 127 measurement days (2 April -- 8 August 2026).

\bigskip
\subsection{Interpreting Geographic Diversity (Q1)}
\label{subsec:5.1.1}

\textbf{Principal finding.} Within the analysed corpus of 250 domains from the Tranco Top 250, macro-regional evaluation reveals that 69.6\% exhibit geographic diversity at the continental scale, resolving to distinct \ac{IP} address pools for probes across at least two continents on one or more measurement days. The rate is identical at national resolution (69.6\%), while 69.2\% of domains also exhibit detectable intra-continental differentiation within the same regional block. Across the 185 countries monitored, domains yield a mean of 13.7 unique \ac{IP} sets, underscoring the fine-grained nature of \ac{CDN} traffic engineering operating well below continental boundaries.

\textbf{Comparison with the literature.}

Van Rijswijk-Deij~\cite{vanRijswijk2016} highlighted that OpenINTEL, operating from a single vantage point in the Netherlands, is structurally blind to any geographic differentiation that authoritative servers implement via anycast or \ac{ECS}. Our distributed architecture --- measurements across 185 countries --- directly addresses this gap. The 69.6\% continent-level diversity rate empirically quantifies the fraction of popular domains whose \ac{DNS} responses OpenINTEL cannot fully characterise: a single-location observatory would observe uniform responses for nearly one third of the Tranco Top-250.

Koch~\cite{Koch2021} showed that the application context determines whether anycast routing introduces operationally significant response variation. The 88.8\% domain-level \ac{NSID} presence rate confirms that anycast-based authoritative \ac{DNS} is the primary structural mechanism underlying geographic response differentiation, consistent with Koch's conclusion.

The mean of 13.7 distinct \ac{IP} address sets at the country level versus 2.65 at the continent level indicates that \ac{CDN} routing granularity is substantially finer than continental boundaries, consistent with Li~\cite{Li2025} and Hours~\cite{Hours2016}: major \ac{CDN} operators maintain country-level routing tables that produce distinct \ac{IP} pools even within the same continent.

\textbf{Inter-campaign finding.} Breaking the analysis by campaign, 60.3\% of common domains are classified as continent-diverse in Campaign~1 (zone-based selection) against 49.7\% in Campaign~2 (continent-stratified fixed set). The lower Campaign~2 rate --- despite broader geographic coverage across Asia-Pacific and Africa --- reflects the reduced EU and NA probe density: fine-grained Western \ac{CDN} routing differences account for a significant share of detectable continent-level diversity, and Campaign~2's 14 European probes are insufficient to capture them. The combined 69.6\% rate benefits from pooling both campaigns over 127 days.

\textbf{Practical implications.}
\begin{itemize}
  \item \textit{For longitudinal \ac{DNS} observatories:} Any measurement system that aims to characterise the \ac{DNS} responses returned to end users in different geographic regions must operate from multiple vantage points. Even a small number of geographically distributed probes suffices to detect geographic variation invisible to single-site measurement.
  \item \textit{For probe selection strategy:} Zone-based selection maximises diversity detection probability in Western \ac{CDN}-heavy corpora; continent-stratified selection reveals richer per-region characterisation at the cost of lower aggregate detection rates. Neither strategy dominates: the optimal design combines a minimum Western probe quota with broad geographic coverage.
  \item \textit{For network performance researchers:} Studies using \ac{DNS} records to infer content-server location should account for geographic diversity; for 69.6\% of popular domains, the \ac{IP} address seen from one location does not represent the \ac{IP} address seen from another.
\end{itemize}

\bigskip
\subsection{Interpreting Temporal Stability (Q2)}
\label{subsec:5.1.2}

\textbf{Principal finding.} The aggregate daily \ac{DNS} response change rate is 38.09\%, and the median per-domain change rate is 48.76\%. Of 250 domains, 17.6\% are fully stable; the remainder exhibit substantial daily churn. The weekly (38.17\%) and monthly (38.18\%) rates are statistically indistinguishable from the daily rate across approximately 97 monthly comparison windows, suggesting that \ac{CDN} pool rotation is a stationary process across all three timescales over four months of continuous measurement.

\textbf{Comparison with the literature.}

Le Pochat~\cite{LePochat2019} measured a 0.6\% daily change rate in the Tranco list composition itself. Our per-domain \ac{DNS} change rates of 38.09\% are therefore approximately 63 times higher than the list-level churn, confirming that a stable corpus of popular domains does not imply stable \ac{DNS} records. The \ac{DNS} layer is substantially more volatile than the popularity-ranking layer, consistent with the short \ac{TTL}s (60--300 seconds) that major \ac{CDN} operators configure to enable rapid traffic redistribution.

Van Rijswijk-Deij~\cite{vanRijswijk2016} demonstrated that \ac{DNS} infrastructure evolves on timescales ranging from hours (\ac{TTL}-driven updates) to years (infrastructure migrations). Our 127-day campaign captures the short-to-medium end of this spectrum. The near-identity of daily, weekly, and monthly change rates (38.09\%, 38.17\%, 38.18\%) confirms that the dominant mechanism is high-frequency \ac{CDN} pool rotation rather than episodic infrastructure events. With approximately 97 monthly comparison windows available, this conclusion rests on substantially stronger empirical ground than a 50-day campaign would provide.

\textbf{Practical implications.}
\begin{itemize}
  \item \textit{Optimal measurement frequency:} For the 17.6\% of fully stable domains, weekly snapshots suffice. For the high-churn majority (median 48.76\%), daily measurement is necessary.
  \item \textit{\ac{DNS} archival strategy:} A tiered policy --- daily collection for volatile domains, weekly for stable ones --- would reduce measurement volume while preserving near-complete change coverage.
\end{itemize}

\bigskip
\subsection{Interpreting the Impact of Geographic Bias (Q3)}
\label{subsec:5.1.3}

\textbf{Principal finding.} Campaign~1's zone-based probe selection concentrates 79.5\% of probe IDs in Europe and North America, consistent with the structural composition of the RIPE Atlas platform; Campaign~2's continent-stratified set reduces this to 20.6\%. Across both campaigns, the Mann-Whitney U test comparing EU+NA against all other continents yields $p = 0.292$, above the $\alpha = 0.05$ threshold: the continental-level imbalance does not statistically distort diversity detection. However, the comparison of 35 single-probe countries against 150 multi-probe countries yields $p \approx 10^{-13}$: countries covered by only one probe observe significantly fewer unique \ac{IP} sets, indicating that within-country probe density matters for the richness of diversity characterisation.

\textbf{Comparison with the literature.}

Bajpai~\cite{Bajpai2017} quantified the RIPE Atlas probe distribution and noted that 91\% of active nodes were confined to Europe and North America. Campaign~1's zone-based selection reproduces this skew (79.5\% EU+NA), confirming that zone rotation alone cannot overcome the structural deficit. Campaign~2 demonstrates that explicit continent-stratified allocation substantially corrects this bias, though the absolute scarcity of RIPE Atlas infrastructure in Africa --- yielding only 39 probes across 34 nations in Campaign~2 --- presents a physical deployment ceiling that algorithmic selection cannot fully overcome.

\textbf{Counter-intuitive finding.} Despite holding only 2.2\% of probe IDs in Campaign~1, Africa detects geographic \ac{DNS} diversity in a share of domains comparable to Europe. This is rooted in \ac{CDN} traffic engineering: operators allocate distinct \ac{IP} pools specifically tailored to Africa, so even a single African probe reliably returns an \ac{IP} footprint segregated from the European pool. The EU+NA imbalance does not create a detection blind spot, but the significant single vs.\ multi-probe result establishes that it does limit the \emph{depth} of characterisation: adding more probes per country reveals additional routing granularity invisible to single-probe coverage.

\textbf{Practical implications.}
\begin{itemize}
  \item \textit{For Q1 robustness:} The 69.6\% geographic diversity rate is robust to the continental probe imbalance (EU+NA $p = 0.292$), since under-represented regions detect comparable per-domain diversity to Europe.
  \item \textit{On the value of stratified allocation:} Campaign~2 demonstrates that explicit continent-stratified allocation substantially corrects the Western probe bias and reveals richer per-region \ac{IP} diversity (e.g., 9.4 vs 5.3 mean unique \ac{IP} sets per domain in Africa). However, the reduction in EU and NA probe density suppresses aggregate diversity detection by approximately 10 percentage points on the common domain corpus. A hybrid design --- continental stratification with a guaranteed minimum EU/NA quota --- would mitigate this trade-off.
  \item \textit{For probe deployment priorities:} Future campaigns should prioritise increasing \emph{per-country} probe density (not just country breadth) in Africa and Asia-Pacific, to transition from binary diversity detection toward statistically robust intra-regional routing characterisation.
\end{itemize}

\bigskip
\subsection{Interpreting Resolver Impact (Q4)}
\label{subsec:5.1.4}

\textbf{Principal finding.} For 52--55\% of (domain, probe) pairs, the three major public resolvers return different \ac{IP} address sets for the same domain queried from the same geographic vantage point. Divergence rates are consistent across all resolver pairs (CF/G: 51.83\%, CF/Q9: 52.81\%, G/Q9: 54.99\%). Cloudflare achieves the lowest median latency (15.5~ms), substantially below Google (24.7~ms) and Quad9 (23.4~ms), reflecting its denser anycast deployment.

\textbf{Comparison with the literature.}

Hours~\cite{Hours2016} demonstrated that clients routed through Google \ac{DNS} were directed to Akamai servers 20~ms farther in \ac{RTT} terms, confirming that resolver choice affects \ac{CDN} edge assignment. Our result extends this finding to a corpus-scale measurement: for more than half of all (domain, probe) pairs across 250 popular domains, the choice between Cloudflare, Google, and Quad9 determines a different \ac{IP} address outcome. This is not a marginal effect confined to specific \ac{CDN} providers --- it is a systematic, corpus-wide property of the resolver ecosystem.

Wang~\cite{Wang2018} documented that routing mismatches caused by resolver centralization can inflate server distance by up to 113\%. Our result provides a measurement-based lower bound: 52--55\% of queries receive a resolver-specific \ac{IP} that differs from at least one other major resolver's answer, implying that the choice of resolver is a first-order determinant of \ac{CDN} edge assignment, comparable in magnitude to the effect of the probe's geographic location established in Q1.

\textbf{Practical implications.}
\begin{itemize}
  \item \textit{For end-user performance:} Users who switch between public resolvers can expect different \ac{CDN} edge assignments for a majority of popular domains, with potential \ac{RTT} and throughput consequences.
  \item \textit{For \ac{DNS} measurement methodology:} Studies that use \ac{DNS} records to infer \ac{CDN} topology or content-server geography must specify which resolver was used; results from Cloudflare, Google, and Quad9 are not interchangeable for more than half of popular domains.
\end{itemize}

\bigskip
\section{Limitations of the Study}
\label{sec:5.2}

\bigskip
\subsection{Methodological Limitations}
\label{subsec:5.2.1}

\textbf{Collection period.} Spanning 127 days (2 April -- 8 August 2026), this campaign captures short-to-medium-term operational dynamics. Broader cyclical or seasonal trends --- strategic \ac{CDN} capacity expansions ahead of peak traffic periods, long-term infrastructural migrations, or the routing impact of geopolitical events --- remain outside the scope of this temporal envelope. Van Rijswijk-Deij~\cite{vanRijswijk2016} demonstrated that multi-year longitudinal observation is necessary to characterise infrastructure evolution at scale; the current results should be understood as a temporally bounded snapshot of an inherently dynamic system.

\textbf{Domain corpus.} Focusing on the Tranco Top-250 limits the scope to the apex of global web popularity. The vast long tail of the \ac{DNS} namespace --- roughly 360 million registered domains~\cite{vanRijswijk2016} --- likely follows entirely different geographic distribution profiles and stability trends. The routing insights documented here reflect the high-traffic domain tier exclusively and cannot be extrapolated to the broader namespace.

\textbf{Geographic probe bias.} While Section~\ref{sec:4.4} establishes that the EU+NA continental imbalance does not prevent diversity detection, the significant single vs.\ multi-probe result confirms that the absolute volume of probes in Africa, South America, and Oceania constrains intra-regional statistical power.

\bigskip
\subsection{Technical Limitations}
\label{subsec:5.2.2}

\textbf{RIPE Atlas platform constraints.} The RIPE Atlas concurrent measurement limit for one-off measurements is a hard practical constraint that must be factored into study design. This limit prevented running Q4 while the 248 ongoing periodic measurements were active, requiring the Q4 campaign to be timed around campaign transitions.

\textbf{\ac{ISP} resolver not available for Q4.} The local \ac{ISP} resolver comparison planned for Q4 could not be executed: the RIPE Atlas \ac{API} rejects the combination of an explicit target \ac{IP} and \texttt{use\_probe\_resolver: true}, and probe inspection showed that a large fraction of probes forward queries to public resolvers anyway, making the condition structurally unimplementable. The Q4 results therefore characterise divergence among the three major public resolvers only, and cannot address whether \ac{ISP} resolvers would produce responses closer to the authoritative signal.

\textbf{Two-campaign heterogeneity.} The two main campaigns used different probe selection strategies (zone-based vs continent-stratified), different probe counts (200 vs 150), and partially different domain corpora (248 vs 201 active domains). Combining them in full-period analyses introduces mixed observational conditions. While the controlled inter-campaign comparison in Section~\ref{subsec:4.1.4} uses the 200 common domains to isolate the strategy effect, full-period results (69.6\% Q1 rate, 38.09\% Q2 rate) pool both conditions and should be interpreted accordingly. Separately, the 9-day pilot phase uses different probe counts and a smaller corpus, reducing its statistical weight in longitudinal analyses.

\bigskip
\subsection{Interpretive Limitations}
\label{subsec:5.2.3}

\textbf{Mechanism attribution.} The \ac{NSID} presence rate difference between diverse and uniform domains is consistent with anycast being the dominant mechanism, but does not allow us to disentangle anycast-based routing from \ac{DNS}-level \ac{CDN} routing (GeoDNS) or \ac{ECS}-driven differentiation, since all three mechanisms can produce different \ac{IP} addresses in different regions.

\textbf{Generalisation beyond the Top-250.} Because the Tranco Top-250 is heavily dominated by enterprise-grade, hyperscale \ac{CDN} deployments, the 69.6\% continent-level diversity rate represents an operational ceiling likely to overestimate the diversity profile of a broader domain sample.

\textbf{Stationarity assumption.} Our conclusion regarding the stationarity of \ac{CDN} pool rotation is supported by approximately 97 monthly comparison windows. While this provides strong evidence for short-to-medium-term stationarity, capturing potential macro-cyclical trends would require extending measurement beyond the current 127-day window.

\bigskip
\section{Contributions of This Work}
\label{sec:5.3}

\subsection{Scientific Contributions}
\label{subsec:5.3.1}

\textbf{Primary empirical contribution.}
This thesis provides a systematic, multi-vantage-point quantification of \ac{DNS} response diversity for the Tranco Top-250, based on 4,146,242 \ac{DNS} measurements from 13,083 unique RIPE Atlas probes across 185 countries over 127 measurement days. The results characterise four dimensions of \ac{DNS} behaviour at a scale and geographic resolution not previously reported for this domain corpus: geographic diversity (Q1), temporal stability (Q2), probe-bias sensitivity (Q3), and resolver-induced divergence (Q4).

\textbf{Empirical validation of literature hypotheses.} Table~\ref{tab:hypothesis_validation} summarises the validation status of key claims from the literature.

\begin{table}[H]
\centering
\small
\begin{tabularx}{\linewidth}{X p{4.0cm} p{5.5cm}}
\toprule
\textbf{Claim} & \textbf{Source} & \textbf{Status} \\
\midrule
Single-vantage-point \ac{DNS} measurement misses geographic variation & ~\cite{vanRijswijk2016} & \textbf{Confirmed}: 69.6\% of domains are continent-diverse \\
Anycast is the primary mechanism of geographic \ac{DNS} variation & ~\cite{Koch2021}; ~\cite{Bortzmeyer2013} & \textbf{Confirmed}: 88.8\% domain-level \ac{NSID} rate \\
\ac{CDN} routing operates at sub-continental granularity & ~\cite{Hours2016}; ~\cite{Li2025} & \textbf{Confirmed}: 13.7 country-level vs 2.65 continent-level \ac{IP} sets \\
\ac{DNS} response content is more volatile than domain popularity rankings & ~\cite{LePochat2019} & \textbf{Confirmed}: 38.09\% daily change rate vs 0.6\% Tranco churn \\
RIPE Atlas geographic bias persists and affects \ac{DNS} measurements & ~\cite{Bajpai2017}; ~\cite{Nosyk2025} & \textbf{Confirmed} (C1: EU+NA 79.5\%; $p=0.292$); single-probe countries observe significantly less diversity ($p \approx 10^{-13}$) \\
Resolver choice affects observed \ac{CDN} edge assignment & ~\cite{Hours2016}; ~\cite{Wang2018} & \textbf{Confirmed}: 52--55\% divergence across all resolver pairs \\
\bottomrule
\end{tabularx}
\caption{Empirical validation status of key claims from the literature}
\label{tab:hypothesis_validation}
\end{table}

\bigskip
\subsection{Methodological Contributions}
\label{subsec:5.3.2}

\textbf{Reproducible measurement pipeline.} The complete pipeline --- from probe selection and RIPE Atlas campaign configuration through data collection, Parquet storage, and statistical analysis --- is documented in Chapter~\ref{ch:method} and published as open-source code on GitHub (\url{https://github.com/ogautier-unam/dns-pipeline}). Key artefacts include the probe selection script with three-round geographic stratification, the daily fetch and parsing pipeline, and the Q1--Q4 analysis scripts.

\textbf{Two-level geographic analysis.} Our two-level analysis (continent + country) yields identical detection rates at both granularities (69.6\% at the continental and national level alike), while \ac{CDN} routing granularity is substantially finer at the country level (mean 13.7 distinct \ac{IP} sets) than at the continent level (mean 2.65 \ac{IP} sets). The separate intra-continental metric (69.2\%) confirms that virtually all continent-diverse domains also exhibit within-continent differentiation. This two-level methodology is reusable for any future distributed \ac{DNS} measurement study.

\textbf{Two-campaign comparative design.} By running consecutive campaigns with zone-based and continent-stratified probe selection, we generate empirical evidence of how probe selection strategy shapes geographic diversity detection. The comparison (Section~\ref{subsec:4.1.4}) shows that zone-based selection detects more aggregate continent-level diversity (60.3\% vs 49.7\% on common domains) while stratified selection reveals richer per-region characterisation. This controlled within-study comparison is a reusable methodological template for future probe selection experiments.

\textbf{RIPE Atlas best practices for \ac{DNS} studies.} We operationalize and build upon the campaign design principles pioneered by Bajpai~\cite{Bajpai2017}, Nosyk~\cite{Nosyk2025}, and Boswell and Perkins~\cite{Boswell2024}: geographic stratification via system tags, direct authoritative-server queries to decouple routing observations from caching artifacts, structured credit budget planning, and \ac{NSID}-based anycast instance identification.

\bigskip
\subsection{Data Contribution}
\label{subsec:5.3.3}

\textbf{Full dataset.} The 127-day measurement campaign (2 April -- 8 August 2026) generated 4,146,242 \ac{DNS} measurement results from 13,083 unique RIPE Atlas probes across 185 countries (Q1--Q3 authoritative measurements). The complete dataset in Parquet format is deposited on Zenodo~\cite{Zenodo} under CC~BY~4.0, making it permanently citable and independently verifiable in accordance with \ac{FAIR} data principles. The pipeline source code is published on GitHub under the MIT licence.

\bigskip
\section{Future Work}
\label{sec:5.4}

\bigskip
\subsection{Immediate Extensions}
\label{subsec:5.4.1}

\textbf{\ac{ISP} resolver comparison for Q4.} The most significant gap left by this study is the absence of a local \ac{ISP} resolver condition in Q4. A dedicated measurement run using \texttt{use\_probe\_resolver: true} --- without a target \ac{IP} field --- would allow comparison of \ac{ISP} resolver responses against the public resolver baselines established here. Prior filtering of probes whose \texttt{/etc/resolv.conf} points to public resolvers (rather than genuine \ac{ISP} infrastructure) would be required to ensure the condition is structurally meaningful.

\textbf{Temporal extension.} Despite covering 127 days, the campaign does not reach the multi-year timescales that Van Rijswijk-Deij~\cite{vanRijswijk2016} identified as necessary for characterising infrastructure evolution. Extending continuous collection would enable detection of seasonal patterns in \ac{CDN} provisioning and longer-cycle routing adjustments. The pipeline is designed for continuous operation and can be restarted with minimal effort.

\bigskip
\subsection{New Research Questions}
\label{subsec:5.4.2}

\textbf{Q5 --- \ac{IPv6} geographic diversity.} Do \ac{DNS} \ac{AAAA} responses exhibit the same geographic diversity as \ac{A} responses? The hypothesis, grounded in Bajpai~\cite{Bajpai2017}, is that \ac{IPv6} \ac{CDN} deployment lags \ac{IPv4}, producing lower inter-regional diversity.

\textbf{Q6 --- \ac{DNS} security and early anomaly detection.} Van der Toorn~\cite{vanderToorn2018} demonstrated that longitudinal \ac{DNS} data enables detection of malicious infrastructure changes 100 days before commercial blacklists. Our dataset, which captures both temporal and spatial dimensions of \ac{DNS} responses, could enable analogous early detection based on anomalous changes in the geographic distribution of \ac{DNS} responses.

\textbf{Q7 --- \ac{DNS} centralisation and digital sovereignty.} The massive proportion of popular domains relying on non-European infrastructure introduces geopolitical frictions regarding digital sovereignty that directly intersect with evolving EU regulatory frameworks such as the NIS2 Directive and the Data Act. Cross-referencing our per-continent diversity metrics with jurisdictional footprints would inject empirical data into this regulatory debate.

\bigskip
\subsection{Methodological Improvements}
\label{subsec:5.4.3}

\textbf{Mechanism attribution.} While our framework flags the presence of geographic \ac{DNS} diversity, it cannot disentangle anycast, GeoDNS, and \ac{ECS} contributions. A custom-configured resolver capable of selectively toggling \ac{ECS} per query would decouple sub-network-driven optimizations from anycast-driven traffic engineering.

\textbf{Probe diversity augmentation.} Campaign~2 demonstrates that continent-stratified allocation can substantially improve Africa and Asia-Pacific coverage (39 probes across 34 African nations vs 239 probes in 37 nations for Campaign~1, a six-fold increase in density). However, 122 of the 136 Campaign~2 countries are still covered by a single probe, leaving them statistically insufficient for intra-regional analysis. Future extensions should set a minimum per-country multi-probe quota (e.g., $\geq$3 probes per country) and integrate complementary platforms such as NLNOG Ring or CAIDA Ark in regions where RIPE Atlas density remains structurally low.

\bigskip
\section{General Conclusion}
\label{sec:5.5}

\subsection{Summary of the Thesis}
\label{subsec:5.5.1}

This thesis addressed the problem of \emph{spatially and temporally distributed \ac{DNS} measurement}: how to design, deploy, and operate a system that captures the geographic diversity and temporal stability of \ac{DNS} responses for the most popular domains on the web, in a manner that is methodologically sound and reproducible.

The motivation was grounded in a well-documented gap: existing large-scale \ac{DNS} observatories~\cite{vanRijswijk2016} operate from single geographic vantage points, making them structurally blind to the geographic differentiation that \ac{CDN} operators~\cite{Koch2021,Hours2016}, \ac{ECS}-aware authoritative servers~\cite{RFC7871,Wang2018}, and \ac{BGP}-anycast routing~\cite{Bortzmeyer2013} deliberately introduce into \ac{DNS} responses. At the same time, the \ac{DNS} resolver ecosystem is increasingly concentrated~\cite{Xu2023}, making resolver choice a confounding variable in any \ac{DNS} measurement study.

Our methodological framework translates these constraints into a tripartite strategy. First, we leverage the RIPE Atlas framework as our distributed measurement infrastructure, deploying a three-round geographic stratification algorithm to systematically isolate probes across 185 countries. Second, to safeguard the integrity of our primary Q1--Q3 campaign, all queries are routed directly to authoritative servers, stripping away recursive-resolver artifacts. Finally, this architecture is augmented by a three-condition public resolver comparison campaign designed to isolate and quantify resolver-induced variations independently.

\bigskip
\subsection{Answers to the Research Questions}
\label{subsec:5.5.2}

\textbf{Q1 --- Geographic diversity:} 69.6\% of the Tranco Top-250 domains return geographically differentiated \ac{DNS} \ac{A} records at the continental scale, with an identical 69.6\% at national resolution (69.2\% show intra-continental differentiation). The mean number of distinct \ac{IP} address pools climbs from 2.65 continentally to 13.7 at the country level. The widespread presence of \ac{NSID}s --- detected across 88.8\% of domains --- points to anycast-driven authoritative \ac{DNS} as the dominant deployment architecture. A traditional, single-point observatory remains blind to this diversity, capturing uniform responses for nearly one third of the web's most popular domains.

\textbf{Q2 --- Temporal stability:} 17.6\% of domains are fully stable over the 127-day campaign. The aggregate daily change rate is 38.09\%; the median per-domain rate is 48.76\%. Daily, weekly, and monthly rates are all nearly identical across $\approx$97 monthly comparison windows, confirming that \ac{CDN} pool rotation is a stationary, high-frequency process. \ac{DNS} response content is approximately 63 times more volatile than the Tranco popularity ranking itself.

\textbf{Q3 --- Probe bias:} Zone-based probe selection (Campaign~1) reproduces the structural RIPE Atlas concentration, with Europe and North America commanding 79.5\% of probe IDs; continent-stratified selection (Campaign~2) reduces this to 20.6\%. In both cases, the EU+NA imbalance does not distort diversity detection ($p = 0.292$), but single-probe countries observe significantly less diversity than multi-probe countries ($p \approx 10^{-13}$), establishing that per-country probe density is a critical determinant of measurement depth.

\textbf{Q4 --- Resolver impact:} For 52--55\% of (domain, probe) pairs, the choice of public resolver determines a different \ac{IP} address outcome. Cloudflare achieves the lowest latency (median 15.5~ms), followed by Google (24.7~ms) and Quad9 (23.4~ms). Resolver selection is thus a first-order determinant of \ac{CDN} edge assignment, comparable in magnitude to the effect of geographic location.

\textbf{Overarching answer.} Distributed, authoritative-querying \ac{DNS} measurement from geographically diverse vantage points is both technically feasible with RIPE Atlas and scientifically necessary to characterise the geographic structure of the modern \ac{DNS}. Systems designed assuming geographic \ac{DNS} uniformity will systematically misrepresent the infrastructure underlying more than two thirds of the most visited websites on the Internet.

\bigskip
\subsection{Lessons Learned}
\label{subsec:5.5.3}

\textbf{Platform lessons.} RIPE Atlas is a powerful distributed measurement platform, but its utility depends on methodological rigour. Key lessons: (i) zone-based probe selection cannot overcome the structural Western concentration of RIPE Atlas --- explicit continent-stratified allocation is required to achieve geographic balance; (ii) however, reducing EU and NA probe density to achieve balance suppresses detection of fine-grained Western \ac{CDN} routing differences --- a minimum Western probe quota must be preserved; (iii) the concurrent measurement limit for one-off measurements is a hard practical constraint on simultaneous campaign scope; (iv) the \ac{ISP} resolver condition is not reliably achievable via the standard RIPE Atlas \ac{API} due to probe-resolver configuration variability.

\textbf{Architectural lessons.} An ingestion pipeline rooted in daily, date-partitioned Parquet files scales seamlessly to handle continuous \ac{DNS} measurement. Our multi-tier storage strategy --- temporary raw JSON under a seven-day retention policy, processed Parquet committed to permanent storage, rclone synchronization to Google Drive --- provides cost-effective disaster recovery on a home-hosted Raspberry Pi with no cloud VM cost.

\textbf{Scientific lessons.} Our data directly refutes the conventional wisdom that geographic \ac{DNS} diversity can be reliably inferred from a single vantage point: 69.6\% of popular domains serve distinct \ac{IP} addresses across different continents. It also establishes that both the client's geographic location and the client's resolver choice are independent, first-order determinants of \ac{DNS} response content.

\bigskip
\subsection{Closing Remarks}
\label{subsec:5.5.4}

The Domain Name System is the invisible substrate on which the Internet's navigability depends. Every web request, email delivery, and \ac{API} call begins with a \ac{DNS} query --- yet the geographic structure of \ac{DNS} responses is rarely measured from the perspectives of the users who depend on it. This thesis is a contribution toward making that structure \emph{visible}, \emph{measurable}, and \emph{reproducible}.

The 4,146,242 measurements collected over 127 days (2 April -- 8 August 2026) constitute a spatiotemporal snapshot of how the Internet's most popular destinations present themselves to users in different parts of the world. They document a \ac{DNS} that is more geographically heterogeneous than single-vantage-point studies suggest (69.6\% of popular domains return different \ac{IP}s in different continents), more temporally volatile than the stability of popularity rankings implies (38.09\% daily change rate), geographically richer to measure with more probes per country (significant single vs.\ multi-probe result), and more resolver-sensitive than commonly assumed (52--55\% inter-resolver divergence).

As Stéphane Bortzmeyer observed~\cite{Bortzmeyer2013}, \ac{DNS} is too often overlooked in studies of Internet resilience and quality of service. With the tools now available --- RIPE Atlas's global probe network, Tranco's manipulation-resistant domain ranking, open Parquet storage, and community-supported analysis infrastructure --- there is no longer a technical barrier to filling this gap. The barrier is methodological care and sustained measurement effort. We hope the dataset, pipeline, and results reported here lower both barriers for the researchers and engineers who come next.
